\documentclass[10pt,a4paper]{amsart}
\usepackage{amsmath,amssymb,amsthm,graphicx}
\usepackage{verbatim}
\newenvironment{reference}{\comment}{\endcomment}
%\newenvironment{reference}{}{}
\newenvironment{slogan}{\comment}{\endcomment}
\newenvironment{history}{\comment}{\endcomment}

% For commutative diagrams we use Xy-pic
\usepackage[all]{xy}

% We use 2cell for 2-commutative diagrams.
\xyoption{2cell}
% Theorem environments.
%
\theoremstyle{plain}
\newtheorem{theorem}[subsection]{Theorem}
\newtheorem{proposition}[subsection]{Proposition}
\newtheorem{lemma}[subsection]{Lemma}

\theoremstyle{definition}
\newtheorem{definition}[subsection]{Definition}
\newtheorem{example}[subsection]{Example}
\newtheorem{exercise}[subsection]{Exercise}
\newtheorem{situation}[subsection]{Situation}

\theoremstyle{remark}
\newtheorem{remark}[subsection]{Remark}
\newtheorem{remarks}[subsection]{Remarks}

\numberwithin{equation}{subsection}

% Macros
%
\def\lim{\mathop{\mathrm{lim}}\nolimits}
\def\colim{\mathop{\mathrm{colim}}\nolimits}
\def\Spec{\mathop{\mathrm{Spec}}}
\def\Hom{\mathop{\mathrm{Hom}}\nolimits}
\def\Ext{\mathop{\mathrm{Ext}}\nolimits}
\def\SheafHom{\mathop{\mathcal{H}\!\mathit{om}}\nolimits}
\def\SheafExt{\mathop{\mathcal{E}\!\mathit{xt}}\nolimits}
\def\Sch{\mathit{Sch}}
\def\Mor{\mathop{\mathrm{Mor}}\nolimits}
\def\Ob{\mathop{\mathrm{Ob}}\nolimits}
\def\Sh{\mathop{\mathit{Sh}}\nolimits}
\def\NL{\mathop{N\!L}\nolimits}
\def\CH{\mathop{\mathrm{CH}}\nolimits}
\def\proetale{{pro\text{-}\acute{e}tale}}
\def\etale{{\acute{e}tale}}
\def\QCoh{\mathit{QCoh}}
\def\Ker{\mathop{\mathrm{Ker}}}
\def\Im{\mathop{\mathrm{Im}}}
\def\Coker{\mathop{\mathrm{Coker}}}
\def\Coim{\mathop{\mathrm{Coim}}}

% Boxtimes
%
\DeclareMathSymbol{\boxtimes}{\mathbin}{AMSa}{"02}

%
% Macros for moduli stacks/spaces
%
\def\QCohstack{\mathcal{QC}\!\mathit{oh}}
\def\Cohstack{\mathcal{C}\!\mathit{oh}}
\def\Spacesstack{\mathcal{S}\!\mathit{paces}}
\def\Quotfunctor{\mathrm{Quot}}
\def\Hilbfunctor{\mathrm{Hilb}}
\def\Curvesstack{\mathcal{C}\!\mathit{urves}}
\def\Polarizedstack{\mathcal{P}\!\mathit{olarized}}
\def\Complexesstack{\mathcal{C}\!\mathit{omplexes}}
% \Pic is the operator that assigns to X its picard group, usage \Pic(X)
% \Picardstack_{X/B} denotes the Picard stack of X over B
% \Picardfunctor_{X/B} denotes the Picard functor of X over B
\def\Pic{\mathop{\mathrm{Pic}}\nolimits}
\def\Picardstack{\mathcal{P}\!\mathit{ic}}
\def\Picardfunctor{\mathrm{Pic}}
\def\Deformationcategory{\mathcal{D}\!\mathit{ef}}
\usepackage{cleveref}
\crefname{originallemma}{original lemma}{original lemmas}
\Crefname{originallemma}{Original lemma}{Original lemmas}
\makeatletter
\newlabel{lemma-exactness-stalks}{{00YL}{source}{}{sourcepointer-1}{}}
\newlabel{lemma-exactness-stalks@cref}{{[originallemma][0][]00YL}{[1][0][]source}{}{}{}}
\makeatother
\crefname{originallemma}{original lemma}{original lemmas}
\Crefname{originallemma}{Original lemma}{Original lemmas}
\makeatletter
\newlabel{lemma-points-recover}{{00Y6}{source}{}{sourcepointer-2}{}}
\newlabel{lemma-points-recover@cref}{{[originallemma][0][]00Y6}{[1][0][]source}{}{}{}}
\makeatother
\crefname{originallemma}{original lemma}{original lemmas}
\Crefname{originallemma}{Original lemma}{Original lemmas}
\makeatletter
\newlabel{lemma-point-pushforward-sheaf}{{00Y8}{source}{}{sourcepointer-3}{}}
\newlabel{lemma-point-pushforward-sheaf@cref}{{[originallemma][0][]00Y8}{[1][0][]source}{}{}{}}
\makeatother
\crefname{originallemma}{original lemma}{original lemmas}
\Crefname{originallemma}{Original lemma}{Original lemmas}
\makeatletter
\newlabel{lemma-mono-epi-sheaves}{{00WN}{source}{}{sourcepointer-4}{}}
\newlabel{lemma-mono-epi-sheaves@cref}{{[originallemma][0][]00WN}{[1][0][]source}{}{}{}}
\makeatother
\crefname{originallemma}{original lemma}{original lemmas}
\Crefname{originallemma}{Original lemma}{Original lemmas}
\makeatletter
\newlabel{categories-lemma-characterize-mono-epi}{{08LR}{source}{}{sourcepointer-5}{}}
\newlabel{categories-lemma-characterize-mono-epi@cref}{{[originallemma][0][]08LR}{[1][0][]source}{}{}{}}
\makeatother
\title{A criterion for a family of points to be conservative}
\author{The Stacks Project Authors}
\date{stacks-project:sites.tex:local-snapshot-20261012; GNU FDL 1.2 or later}
\begin{document}
\maketitle
\noindent\textit{Editorial note.} All delivered mathematical text is the author's own, copied from the pinned Stacks Project chapter file sites-cohomology.tex; no statement or proof step was rewritten, reordered or removed.\par\medskip
\noindent\textit{Editorial note.} The proof is the author's complete argument; every step the author writes out is inside the counted span and no closing sentence is a placeholder.\par\medskip
\noindent\textit{Editorial note.} Every labelled reference occurring inside the counted spans is carried in assembly.external\_references as a non-load-bearing navigation pointer (same-chapter ones at their definition site, cross-chapter ones at the invocation inside the counted span), so no load-bearing mathematics is replaced by a stub. No bibliography citation occurs inside the counted spans.\par\medskip
\noindent\textit{Editorial note.} The checked chapter list for same-name collisions is the set of stacks chapters staged in this repository: chow, cohomology, divisors, etale-cohomology, flat, more-algebra, more-morphisms, morphisms, spaces-more-morphisms, varieties together with this card's chapter sites-cohomology.\par\medskip
\noindent\textit{Editorial note.} The preamble is the chapter's own retained minimal preamble (sources/stacks-sites-cohomology/retained-preamble.tex), a verbatim excerpt of the native preamble: the Xy-pic setup and the native theorem declarations with the mathematical macros.\par\medskip
\noindent\textit{Allowed standard prerequisites (editorial).}\begin{itemize}
\item A conservative family detects equality of maps. Same chapter, Lemma lemma-exactness-stalks. Supplies the detection of inequality of the maps induced by two distinct sections.
\item Stalks of representable sheaves. Same chapter, Lemma lemma-points-recover and Lemma lemma-point-pushforward-sheaf. Identify the stalk of the representable sheaf with the value of the functor.
\item Injections and surjections as morphisms criteria. Same chapter, Lemma lemma-mono-epi-sheaves and Categories, Lemma categories-lemma-characterize-mono-epi. Allow injectivity and surjectivity to be checked through the diagonal and the codiagonal.
\end{itemize}\medskip
\ifdefined\hypertarget\hypertarget{sourcepointer-1}{}\fi
\noindent\textit{Source reference.} Original lemma 00YL: sources/stacks-sites/pinned-original-source.tex, source lines 9324--9324. Navigation pointer invoked at proof line 9324; the counted span contains the pointer only and the cited result is not reproduced inside it, so the pointer is registered as a non-load-bearing navigation stub in the established form for Stacks material. No bibliography citation is involved.\par\medskip
\ifdefined\hypertarget\hypertarget{sourcepointer-2}{}\fi
\noindent\textit{Source reference.} Original lemma 00Y6: sources/stacks-sites/pinned-original-source.tex, source lines 9326--9326. Navigation pointer invoked at proof line 9326; the counted span contains the pointer only and the cited result is not reproduced inside it, so the pointer is registered as a non-load-bearing navigation stub in the established form for Stacks material. No bibliography citation is involved.\par\medskip
\ifdefined\hypertarget\hypertarget{sourcepointer-3}{}\fi
\noindent\textit{Source reference.} Original lemma 00Y8: sources/stacks-sites/pinned-original-source.tex, source lines 9327--9327. Navigation pointer invoked at proof line 9327; the counted span contains the pointer only and the cited result is not reproduced inside it, so the pointer is registered as a non-load-bearing navigation stub in the established form for Stacks material. No bibliography citation is involved.\par\medskip
\ifdefined\hypertarget\hypertarget{sourcepointer-4}{}\fi
\noindent\textit{Source reference.} Original lemma 00WN: sources/stacks-sites/pinned-original-source.tex, source lines 9338--9338. Navigation pointer invoked at proof line 9338; the counted span contains the pointer only and the cited result is not reproduced inside it, so the pointer is registered as a non-load-bearing navigation stub in the established form for Stacks material. No bibliography citation is involved.\par\medskip
\ifdefined\hypertarget\hypertarget{sourcepointer-5}{}\fi
\noindent\textit{Source reference.} Original lemma 08LR: sources/stacks-sites/pinned-original-source.tex, source lines 9343--9343. Navigation pointer invoked at proof line 9343; the counted span contains the pointer only and the cited result is not reproduced inside it, so the pointer is registered as a non-load-bearing navigation stub in the established form for Stacks material. No bibliography citation is involved.\par\medskip
\begin{lemma}
\label{lemma-enough}
Let $\mathcal{C}$ be a site and let $\{(p_i, u_i)\}_{i\in I}$ be a
family of points. The family is conservative if and only if for every
sheaf $\mathcal{F}$ and every $U\in \Ob(\mathcal{C})$ and every
pair of distinct sections $s, s' \in \mathcal{F}(U)$, $s \not = s'$ there
exists an $i$ and $x\in u_i(U)$ such that the triples
$(U, x, s)$ and $(U, x, s')$ define distinct elements of
$\mathcal{F}_{p_i}$.
\end{lemma}

\begin{proof}
Suppose that the family is conservative and that $\mathcal{F}$, $U$, and
$s, s'$ are as in the lemma. The sections $s$, $s'$ define maps
$a, a' : (h_U)^\# \to \mathcal{F}$ which are distinct. Hence, by Lemma
\ref{lemma-exactness-stalks} there is an $i$ such that $a_{p_i}
\not = a'_{p_i}$. Recall that $(h_U)^\#_{p_i} = u_i(U)$, by
Lemmas \ref{lemma-points-recover} and
\ref{lemma-point-pushforward-sheaf}.
Hence there exists an $x \in u_i(U)$ such that $a_{p_i}(x)
\not = a'_{p_i}(x)$ in $\mathcal{F}_{p_i}$.
Unwinding the definitions you see that $(U, x, s)$
and $(U, x, s')$ are as in the statement of the lemma.

\medskip\noindent
To prove the converse, assume the condition on the existence of
points of the lemma. Let $\phi : \mathcal{F} \to \mathcal{G}$
be a map of sheaves which is an isomorphism at all the stalks.
We have to show that $\phi$ is both
injective and surjective, see Lemma \ref{lemma-mono-epi-sheaves}.
Injectivity is an immediate consequence of the assumption.
To show surjectivity we have to show that
$\mathcal{G} \amalg_\mathcal{F} \mathcal{G} \to \mathcal{G}$
is an isomorphism
(Categories, Lemma \ref{categories-lemma-characterize-mono-epi}).
Since this map is clearly surjective, it suffices to check injectivity
which follows as $\mathcal{G} \amalg_\mathcal{F} \mathcal{G} \to \mathcal{G}$
is injective on all stalks by assumption.
\end{proof}

\end{document}
